\section{The cubic circuit cap}\label{sec:cap}
Put $a=\essinf G$. Choose $M\ge1$ so large that $q=G((M,\infty))<1/1024$.
Choose $a<a_o<b_o<M$ with $G([0,a_o])>0$ and $G([0,b_o))<1/1024$.
Nonatomicity at $a$ permits this choice even for singular continuous $G$.
Unbounded support permits $M<a_c<b_c<\infty$ with
$G((M,a_c])>0$ and $G([b_c,\infty))>0$. Set
\begin{align*}
 L_o&=[0,a_o],&H_o&=[b_o,M],&r_o&=G([0,M]),\\
 L_c&=(M,a_c],&H_c&=[b_c,\infty),&r_c&=G((M,\infty)),\\
 p_{a,i}&=G(L_i),&p_{b,i}&=G(H_i),&
 g&=\min(b_o-a_o,b_c-a_c)>0,\\
 p_*&=\min_{i\in\{o,c\}}p_{a,i}/r_i>0.
\end{align*}
All four band masses are positive. Call an original edge \emph{usable} when
$t_e\in H_o\cup H_c$. Its bad probability $\epsilon<1/512$ satisfies
$8\sqrt\epsilon<1$ and $6\sqrt\epsilon<1/2$. These bands remain fixed.

Let $V_e$ be independent uniform variables on $(0,1)$, independent of $t$;
write $X=(t,V)$. An edge is open when $t_e\le M$. Let $L_e$ be the minimum
length of an open simple circuit strictly enclosing both endpoints of $e$,
and define
\begin{equation}\label{eq:cap}
 K_e=M(L_e^3+V_e),\qquad s_e=\min(t_e,K_e),\qquad T=T_s.
\end{equation}
The definition uses the minimum length only, not a choice among tied circuits.

\begin{lemma}\label{lem:capfacts}
Almost surely all $L_e$ exist, $L_e\ge10$, and
$\PP(L_e>r)\le C\exp(-c\sqrt r)$. The field $s$ is a square-symmetric
factor of iid, mixing under every nonzero lattice translation, has every
finite moment, and has unique finite geodesics. Every open edge is unchanged.
\end{lemma}
\begin{proof}
Failure of an open circuit in a suitable annulus of outer radius $R$ forces
a closed dual crossing of length at least $cR$. Erasing loops and counting
at most three choices after the first edge bounds this probability by
$CR\sum_{j\ge cR}(3q)^j$. A simple circuit in the outer box has $O(R^2)$
edges. This gives the stated tail and, by countability, existence everywhere.
The same argument surrounds every finite box by open circuits arbitrarily
far away. A lattice circuit strictly enclosing a horizontal edge has
horizontal span at least three and vertical span at least two; hence its
length is at least ten. The vertical case is identical. In particular
$K_e\ge2ML_e>M$, so open edges are unchanged and
$0<s_e\le M(L_e^3+1)$ has all moments.

Conditional on all open indicators, the original magnitudes have independent
nonatomic conditional laws, the $V_e$ remain independent uniforms, and the
$L_e$ are fixed. Thus the $s_e$ are conditionally independent and nonatomic.
For two distinct simple paths, condition on an edge in their symmetric
difference to rule out equal sums; countability treats all finite paths.
One must not freeze that edge's $V_e$ in this no-tie argument. Existence of
finite minimizers is established by the coercivity argument in
Lemma~\ref{lem:metric}. Covariance and symmetries of the minimum-length
construction give the factor and mixing assertions.
\end{proof}

Choose any measurable minimum circuit $C_e$ for estimates, and put
\[E_x=\{e:x\text{ is strictly inside }C_e\}.\] The tail above and the fact
that $x$ inside $C_e$ implies $\operatorname{dist}(x,e)\le C L_e$ show
that $E_x$ is finite, simultaneously for every $x$.

\begin{proposition}[Exact preservation]\label{prop:equivalence}
On one event, the sets of geometric $t$- and $s$-bigeodesics are exactly
equal. Every edge of each such path is unchanged. Every $s$-ray from a fixed
root $x$ is uncapped outside the finite set $E_x$, simultaneously over all
rays from that root.
\end{proposition}
\begin{proof}
Call $e$ changed if $s_e<t_e$, so $s_e=K_e$. A bigeodesic in either metric
using $e$ must enter and leave $C_e$. Its segment between contacts bracketing
$e$ costs at least $K_e$, whereas a shorter boundary arc costs at most
$ML_e/2$. Thus neither kind of bigeodesic uses a changed edge.

For the nontrivial implication, let $\gamma$ be a $t$-bigeodesic and let
$\pi$ be a finite $s$-walk between $a,b\in\gamma$. Eliminate a changed
occurrence $e$ of $\pi$. If both endpoints lie outside or on $C_e$, replace
the bracketed crossing by a boundary arc. If just $a$ lies inside, replace
the prefix through $e$ by a segment of $\gamma$ from $a$ to a boundary
contact and then a boundary arc to a later contact of $\pi$; reverse the
construction if just $b$ is inside. The part of $\gamma$ between its
outermost contacts with $C_e$ costs at most $ML_e/2$, by $t$-minimality,
so the inserted portion costs at most $ML_e$. If both endpoints are inside,
replace all of $\pi$ by $\gamma[a,b]$, contained in that same outermost-contact
interval and costing at most $ML_e/2$.

Every removed portion contains an edge of $s$-cost $K_e\ge2ML_e$, and every
inserted edge lies on the unchanged circuit or the unchanged $\gamma$.
The number of changed occurrences strictly decreases, with no recursive
creation of changed edges. After finitely many steps the walk is entirely
unchanged and costs no more than the original $s$-walk. Its $t$-cost equals
its $s$-cost, so $t$-minimality of $\gamma$ proves $s$-minimality.
Conversely an unchanged $s$-geodesic is $t$-minimal because $s\le t$.
The proof is pathwise and takes no union over an uncountable path space.

If a ray from $x$ uses changed $e\notin E_x$, it begins outside or on $C_e$
and eventually exits that finite region. The same boundary shortcut is a
contradiction. This proves the rooted-ray statement without asserting any
reattachment of a preserved tail to the same root in the other metric.
\end{proof}

\begin{lemma}[Finite influence and monotonicity]\label{lem:influence}
On one original full event, any finite modification of the positive finite
original weights, with $V$ fixed, changes only finitely many caps. If the
original weights increase, then $s'\ge s$. Under such an upward modification,
an uncapped path avoiding the modified original set is unchanged.
\end{lemma}
\begin{proof}
Tail summability gives $L_e/(1+|e|)\to0$ uniformly over distant edges.
Circuits around every finite box ensure that protecting circuits still exist
after any finite modification $F$. For distant $e$, every circuit enclosing
$e$ and meeting $F$ is longer than $L_e$. An old minimizer avoids $F$, so
$L'_e\le L_e$; a new minimizer of that length also avoids $F$, so it was
old-admissible and $L_e\le L'_e$. Thus equality holds far away.
Raising originals only removes open circuits, giving $L'_e\ge L_e$,
$K'_e\ge K_e$, and $s'_e\ge s_e$. In this upward case an uncapped edge
whose original weight is not modified still has $s'_e=t_e=s_e$.
\end{proof}

\begin{proposition}[Failure of upward finite energy]\label{prop:energy}
For unbounded $G$, the cap law fails DH's upward finite-energy condition
and AH A2(v).
\end{proposition}
\begin{proof}
For $f\ne e$, its status is visible from $s_f$: $s_f\le M$ exactly when
$t_f\le M$. A circuit strictly enclosing both endpoints of $e$ cannot
contain $e$, so $L_e$ is determined by the outside capped field $s_{-e}$.
Consequently $s_e\le M(L_e^3+1)$ conditionally on $s_{-e}$. The event
$L_e=10$ has positive probability and is outside-measurable. On this event
the conditional probability of $s_e\ge\lambda:=1001M+1$ is zero.
Unconditionally that probability is positive: a finite closed pattern
can exclude all short surrounding circuits, an open outer circuit has
positive probability, and independently $t_e>\lambda$ has positive
probability. Take the excluded length sufficiently large that its cap
exceeds $\lambda$. Thus the positive unconditional tail is not conditionally
available everywhere.
\end{proof}
No iid law for $s$ or restoration of its conditional finite energy is used
later. Optional endpoint corrections $T_t-T_s\le B(x)+B(y)$ with
$B(x)=\sum_{e\in E_x}t_e<\infty$ do not supply uniform original shape and
are not needed in the essential chain.
